\documentclass[11pt]{article}
\usepackage[margin=1in]{geometry}
\usepackage{amsmath}
\usepackage{amssymb}
\usepackage{booktabs}
\usepackage{longtable}
\usepackage{array}
\usepackage{graphicx}
\usepackage{float}
\usepackage{hyperref}
\usepackage{setspace}

\hypersetup{
    colorlinks=true,
    linkcolor=black,
    citecolor=black,
    urlcolor=blue
}

\setstretch{1.15}

\title{Machine Learning for Dip-Coated Thin-Film Thickness Prediction:\\
From Classical Coating Theory to Bayesian-Optimized Random Forests}
\author{Pitt ML MSE Project\\
\small Replace with author names and affiliations before submission}
\date{\today}

\begin{document}
\maketitle

\begin{abstract}
Dip coating is one of the most widely used wet-processing routes for depositing functional thin films, yet the predictive modeling of final film thickness remains difficult when the process involves multivariable material chemistry, nonlinear drying or bonding effects, and limited experimental data. Classical dip-coating theory, especially the Landau--Levich framework, provides a powerful first-order relation between withdrawal conditions and entrained liquid-film thickness, but it does not directly model post-deposition bonded thickness in multicomponent processing systems. This work documents the complete evolution of the Pitt ML MSE project from a small experimental dataset through synthetic data augmentation, holdout validation, full-dataset cross-validation, ensemble modeling, and Bayesian hyperparameter optimization. The project began with a 58-sample experimental dataset containing concentration, uncoated layer thickness, total thickness, and bonded thickness. Two synthetic-data stages were then used: a 1000-sample correlation-preserving dataset for preliminary model development and a 3000-sample feature-expanded dataset that added polarity, viscosity, boiling point, and surface tension to support a more realistic validation study.

Five supervised regression families were investigated: linear regression, $k$-nearest neighbors, backpropagation neural networks, random forest regression, and gradient boosting regression. On the 1000-sample random split study, the optimized random forest achieved the best test performance with $R^2 = 0.8082$. A more difficult sequential holdout experiment on the 3000-sample dataset exposed severe train--validation distribution shift and initially drove all baseline models to negative $R^2$ values; after feature expansion from 3 to 7 variables and model-specific optimization, the holdout random forest recovered to $R^2 = 0.3284$. To remove the artifacts of sequential partitioning, stratified 5-fold cross-validation based on target quantile bins was then applied to the full dataset. Under this more robust protocol, the optimized random forest reached an out-of-fold $R^2 = 0.9716$, a stacked random forest plus gradient boosting ensemble reached $R^2 = 0.9722$, and a final Bayesian hyperparameter search over the full dataset produced a tuned random forest with out-of-fold $R^2 = 0.9734$, outperforming all other tested approaches.

Beyond prediction accuracy, this study shows how validation design strongly shapes conclusions in small-data materials workflows. The results demonstrate that the principal challenge was not simply model choice, but the interaction of feature completeness, synthetic data strategy, and split design. The final optimized random forest provides the strongest predictive ceiling in this project, while residual and feature-importance analysis identifies clear next steps in physicochemical interaction engineering and constraint-aware synthetic data generation.
\end{abstract}

\section{Introduction}

\subsection{Dip coating as a foundational thin-film process}
Dip coating is a standard manufacturing and laboratory technique for depositing thin films on solid substrates. In a typical experiment, a substrate is immersed in a liquid bath and subsequently withdrawn at controlled speed. During withdrawal, viscous drag entrains liquid upward while gravity and surface tension oppose film thickening. The retained liquid layer later evolves through drainage, evaporation, solute redistribution, and in some cases polymer consolidation or bonding. Dip coating is therefore relevant across protective coatings, electronics, optics, biomaterials, and polymer processing, where final film quality depends on thickness uniformity, repeatability, and sensitivity to coupled processing variables.

From a fluid-mechanics perspective, the starting point for many dip-coating analyses is the Landau--Levich law, later discussed within the broader Landau--Levich--Derjaguin framework. For an ideal Newtonian liquid on a smooth substrate, the wet film thickness scales as
\begin{equation}
h \approx 0.94 \, \ell_c \, \mathrm{Ca}^{2/3},
\end{equation}
where $\ell_c = \sqrt{\gamma/(\rho g)}$ is the capillary length and $\mathrm{Ca} = \mu U/\gamma$ is the capillary number defined through viscosity $\mu$, withdrawal velocity $U$, and surface tension $\gamma$ \cite{landau1942,rio2017}. This relation is elegant because it compresses the coating problem into a compact scaling law that balances viscous entrainment and capillary restoration.

\subsection{Why classical models are not enough}
Although the Landau--Levich family of models is physically indispensable, it is also structurally limited for the present problem. First, it predicts the thickness of an entrained liquid layer, not the final bonded thickness after downstream physical and chemical evolution. Second, the classical derivation assumes Newtonian behavior, ideal interfaces, and relatively simple substrate conditions. Third, real dip-coating systems are often influenced by interacting variables that do not appear naturally in the canonical formula, including concentration, solvent polarity, viscosity, boiling point, and substrate-adjacent thickness states. Finally, in practical laboratory settings the number of directly measured experiments is often too small to support a purely mechanistic closure.

These limitations are visible in the present project. The original experimental dataset contained only 58 observations, and the target was bonded thickness rather than liquid entrainment thickness. The process therefore sits precisely in the regime where physics provides valuable intuition but not a complete predictive surrogate. For such problems, machine learning offers an attractive alternative: it can learn nonlinear interactions among process variables, incorporate heterogeneous descriptors, and produce high-fidelity surrogates once enough representative training data are available \cite{schmidt2019,kailkhura2019}.

\subsection{Machine learning as a new avenue for dip-coating prediction}
The central premise of this project is that bonded thickness can be modeled more effectively by combining physically meaningful input variables with nonlinear regression methods. Rather than replacing coating physics, machine learning is used here as a complementary surrogate framework. Process variables and material descriptors encode the state of the system; regression models infer the nonlinear mapping to bonded thickness; and careful validation determines whether the learned relation is robust or merely overfit.

The contribution of this work is not a single model benchmark in isolation. Instead, it presents the full research arc of the project:
\begin{enumerate}
    \item establishing baseline machine-learning models on a small experimental-to-synthetic workflow,
    \item diagnosing failure under sequential holdout validation,
    \item correcting that failure through feature expansion and better validation design,
    \item evaluating ensemble methods for incremental gains, and
    \item performing Bayesian hyperparameter optimization on the strongest model to identify ceiling performance.
\end{enumerate}

The final outcome is a detailed and reproducible account of how predictive performance evolved from modest random-split results to a high-performing full-dataset model, and why that evolution occurred.

\section{Background and Related Work}

\subsection{Classical dip-coating theory}
The classical theory of plate withdrawal developed by Landau and Levich remains one of the cornerstones of coating science \cite{landau1942}. Later reviews have shown that the associated scaling relations remain highly effective for Newtonian dip coating on ideal substrates, while also emphasizing that deviations emerge when interfacial rheology, non-Newtonian bulk behavior, evaporation, particulate loading, or substrate complexity become important \cite{rio2017}. These deviations matter because real engineering coating systems are rarely ideal. As a result, even when classical scaling captures qualitative trends, it may fail to predict the final solid-state film with sufficient fidelity for materials design.

\subsection{Machine learning in materials and process modeling}
Machine learning has become increasingly important in materials science because it can identify predictive patterns in high-dimensional experimental or simulated data that would be difficult to express analytically \cite{schmidt2019,kailkhura2019}. This is especially useful when data include mixed feature types, nonlinear couplings, and latent process interactions. In the present work, the feature set spans direct process variables such as concentration and thickness measurements as well as physicochemical solvent descriptors such as polarity, viscosity, boiling point, and surface tension. Such mixed descriptors are well suited to supervised learning methods, especially tree ensembles that can capture interaction structure without requiring explicit polynomial expansion.

\subsection{Model families relevant to this project}
The project investigated five standard supervised regression families:
\begin{itemize}
    \item \textbf{Linear regression}, used as a transparent baseline for approximately linear response surfaces.
    \item \textbf{$k$-nearest neighbors (KNN)}, a local nonparametric method whose predictions depend on nearby training instances \cite{aha1991}.
    \item \textbf{Backpropagation neural networks (BPNN)}, multilayer feed-forward networks trained by gradient-based weight updates \cite{rumelhart1986}.
    \item \textbf{Random forest regression (RFR)}, an ensemble of decorrelated decision trees whose averaged output often provides strong nonlinear performance and natural robustness to feature scaling \cite{breiman2001}.
    \item \textbf{Gradient boosting regression (GBR)}, a stagewise additive ensemble of weak learners optimized to correct prior residuals \cite{friedman2001}.
\end{itemize}

These model families were chosen because together they span linear, local, neural, and tree-based inductive biases. Their differing behavior under feature incompleteness and distribution shift is one of the main interpretive threads of the project.

\subsection{Hyperparameter optimization}
Modern machine-learning performance depends strongly on hyperparameter choice. Exhaustive grid search is straightforward but often inefficient when the space is large. Bayesian optimization instead treats validation performance as a black-box objective and adaptively samples promising parameter configurations based on prior evaluations \cite{snoek2012}. In this project, Bayesian optimization was applied to the strongest single model family, random forest regression, during the final full-dataset phase to determine whether a higher predictive ceiling could be reached beyond manual tuning and ensemble blending.

\section{Project Scope and Research Questions}
This paper addresses five research questions:
\begin{enumerate}
    \item Can machine-learning models predict bonded thickness from dip-coating process variables and solvent descriptors more effectively than simple linear baselines?
    \item How sensitive are model conclusions to dataset construction and validation protocol?
    \item Which model family is most robust under severe train--validation distribution shift?
    \item Do random forest and gradient boosting ensembles offer a meaningful improvement over a tuned random forest?
    \item Does Bayesian optimization of the best-performing model improve ceiling performance on the full dataset?
\end{enumerate}

\section{Data and Experimental Design}

\subsection{Original experimental dataset}
The project began with a primary aggregated experimental dataset stored in \texttt{data/agg.data.xlsx}. It contains 58 rows and 4 columns:
\begin{enumerate}
    \item Concentration (g/mL)
    \item Uncoated Layer (nm)
    \item Total Thickness (nm)
    \item Bonded Thickness (nm)
\end{enumerate}

The learning target throughout the project was bonded thickness. This is important because bonded thickness is not simply the wet entrained film thickness predicted by classical dip-coating scaling. It is a downstream response influenced by both deposition and post-deposition evolution.

\subsection{Synthetic data generation stages}
Because the original dataset was too small for broad model comparison, two synthetic-data stages were introduced.

\paragraph{Stage 1: CTGAN-generated synthetic data.}
An early synthetic-data workflow used a conditional tabular generative adversarial network (CTGAN) through the SDV implementation. This stage produced \texttt{data/synthetic\_data.csv}, intended to expand the training pool while preserving basic tabular structure. CTGAN is designed for mixed-type tabular synthesis and is widely used when direct experimental data are limited \cite{xu2019}.

\paragraph{Stage 2: correlation-preserving improved synthetic data.}
The project then moved to a more controlled generator, \texttt{src/generate\_synthetic\_data\_improved.py}, which bootstraps observations from the original dataset and adds correlated noise derived from the empirical correlation matrix. This produced \texttt{data/synthetic\_data\_improved.csv}, a 1000-sample dataset with the same 4 columns as the experimental source. The goal of this stage was to preserve feature--target correlation structure more faithfully than unconstrained generative synthesis.

\paragraph{Stage 3: feature-expanded 3000-sample holdout dataset.}
The final large-scale study used \texttt{Holdout Validation/Data/synthetic\_data\_improved.csv}, a 3000-sample dataset with 8 columns:
\begin{enumerate}
    \item Polarity (XLogP3)
    \item Viscosity (cP)
    \item Boiling Point (K)
    \item Surface Tension (mN/m)
    \item Concentration (g/mL)
    \item Uncoated Layer (nm)
    \item Total Thickness (nm)
    \item Bonded Thickness (nm)
\end{enumerate}

This dataset enabled a richer study in which chemistry-aware descriptors could complement the original three process/thickness variables.

\subsection{Dataset summary}
Table~\ref{tab:datasets} summarizes the three core datasets used in the project.

\begin{table}[H]
\centering
\caption{Datasets used throughout the project.}
\label{tab:datasets}
\begin{tabular}{p{3.4cm}p{2.0cm}p{1.6cm}p{7.0cm}}
\toprule
Dataset & Samples & Features & Role in project \\
\midrule
\texttt{agg.data.xlsx} & 58 & 3 inputs + 1 target & Original experimental dataset used as the physical starting point for all downstream work. \\
\texttt{data/synthetic\_data\_improved.csv} & 1000 & 3 inputs + 1 target & Preliminary model development and optimization under random train/test splitting. \\
\texttt{Holdout Validation/Data/synthetic\_data\_improved.csv} & 3000 & 7 inputs + 1 target & Sequential holdout study, stratified full-dataset cross-validation, ensemble analysis, and Bayesian optimization. \\
\bottomrule
\end{tabular}
\end{table}

\subsection{Observed statistical characteristics of the 3000-sample dataset}
The 3000-sample feature-expanded dataset had mean bonded thickness of approximately 0.525 nm with standard deviation 0.183 nm. Concentration and total thickness were the strongest positively correlated variables in the full dataset, while polarity and viscosity showed weak negative correlations with the target. Importantly, the synthetic generator allowed limited extrapolation beyond the empirical min/max of the original data. As a result, the expanded dataset contains small negative values in concentration and total thickness, which are not strictly physical but were accepted in the project as part of the synthetic stress-testing framework. This issue is revisited in the limitations section.

\section{Modeling Methodology}

\subsection{Prediction target and notation}
Let $x \in \mathbb{R}^p$ denote the feature vector and let $y$ denote bonded thickness. The learning task is supervised regression: estimate a function $\hat{f}(x)$ such that
\begin{equation}
\hat{y} = \hat{f}(x)
\end{equation}
approximates the underlying process mapping from material/process descriptors to bonded thickness.

\subsection{Evaluation metrics}
Model performance was measured with mean squared error (MSE), root mean squared error (RMSE), mean absolute error (MAE), and coefficient of determination ($R^2$):
\begin{align}
\mathrm{MSE} &= \frac{1}{n}\sum_{i=1}^{n}(y_i-\hat{y}_i)^2, \\
\mathrm{RMSE} &= \sqrt{\mathrm{MSE}}, \\
\mathrm{MAE} &= \frac{1}{n}\sum_{i=1}^{n}|y_i-\hat{y}_i|, \\
R^2 &= 1 - \frac{\sum_{i=1}^{n}(y_i-\hat{y}_i)^2}{\sum_{i=1}^{n}(y_i-\bar{y})^2}.
\end{align}

\subsection{Linear regression}
Linear regression assumes an affine relationship:
\begin{equation}
\hat{y} = \beta_0 + \sum_{j=1}^{p}\beta_j x_j.
\end{equation}
Its main role in this project was interpretive: if linear regression performed comparably to nonlinear methods, then the additional complexity of tree or neural models would not be justified.

\subsection{K-nearest neighbors regression}
KNN regression predicts from nearby training instances:
\begin{equation}
\hat{y}(x) = \frac{\sum_{i \in \mathcal{N}_k(x)} w_i y_i}{\sum_{i \in \mathcal{N}_k(x)} w_i},
\end{equation}
where $\mathcal{N}_k(x)$ is the set of the $k$ closest observations to $x$ and $w_i$ is either uniform or distance-based \cite{aha1991}. KNN is flexible but can be fragile under distribution shift because it relies directly on neighborhood geometry.

\subsection{Backpropagation neural network}
The BPNN used in the project is a multilayer feed-forward neural network. In compact notation, a depth-$L$ network computes
\begin{equation}
\hat{y} = W_L \sigma\left(W_{L-1}\sigma\left(\cdots \sigma(W_1 x + b_1)\right)+b_{L-1}\right)+b_L,
\end{equation}
with parameters optimized through backpropagation \cite{rumelhart1986}. BPNNs can capture nonlinear structure but generally require careful regularization and enough representative data to avoid unstable generalization.

\subsection{Random forest regression}
Random forest regression averages predictions from many decorrelated trees:
\begin{equation}
\hat{y}(x) = \frac{1}{T}\sum_{t=1}^{T} f_t(x),
\end{equation}
where each tree $f_t$ is trained on a bootstrap sample and split using a random subset of features \cite{breiman2001}. The project repeatedly found that random forests were the most robust family under incomplete feature sets and split-induced distribution shift.

\subsection{Gradient boosting regression}
Gradient boosting constructs an additive model in stages:
\begin{equation}
F_M(x) = \sum_{m=1}^{M}\nu \gamma_m h_m(x),
\end{equation}
where each weak learner $h_m$ is fitted to improve the residuals from earlier stages \cite{friedman2001}. Gradient boosting is often highly competitive but can be more sensitive than random forests to hyperparameter choices and data pathologies.

\subsection{Model training chronology}
The methodology evolved through four phases:
\begin{enumerate}
    \item \textbf{Preliminary model comparison on the 1000-sample improved synthetic dataset}: all four main model families were trained under random 80/20 train/test splits.
    \item \textbf{Sequential holdout validation on the 3000-sample feature-expanded dataset}: the first 2000 rows were used for training and the last 1000 for validation.
    \item \textbf{Stratified full-dataset cross-validation}: target-quantile bins were used to create approximately balanced folds for regression validation on the full 3000-sample dataset.
    \item \textbf{Ensemble analysis and Bayesian optimization}: the strongest models from the previous stages were combined or tuned to identify ceiling performance.
\end{enumerate}

\subsection{Feature scaling policy}
Feature scaling was used for neural-network and distance-based models, following the project implementation. Tree-based models were ultimately evaluated without relying on scaling because their split-based structure is not inherently scale dependent.

\subsection{Sequential holdout design}
The holdout-validation study used a deterministic split:
\begin{itemize}
    \item training set: rows 0--1999,
    \item validation set: rows 2000--2999.
\end{itemize}
This design was intentionally retained during diagnosis because it exposed a major project issue: the first two-thirds and last third of the dataset did not come from the same empirical distribution.

\subsection{Stratified $k$-fold design for regression}
To overcome the shortcomings of sequential splitting, the full-dataset validation phase used 5-fold cross-validation with target quantile bins. Continuous bonded-thickness values were discretized into 10 quantile strata, and folds were created with balanced target-range representation. This is not textbook stratification for classification; rather, it is a practical regression-oriented surrogate to keep each fold comparable in response distribution. The same fold design was reused for full-dataset random forest evaluation, ensemble comparison, and Bayesian optimization to maintain comparability.

\subsection{Ensemble methodology}
Two ensemble strategies were tested on the full dataset:
\begin{enumerate}
    \item a fixed 50/50 average of random forest and gradient boosting predictions,
    \item an adaptive blend, where the random forest weight was chosen by inner cross-validation over a grid from 0.0 to 1.0 in steps of 0.1.
\end{enumerate}
In addition, a stacked ensemble with a ridge-regression meta-learner was trained on out-of-fold base predictions from random forest and gradient boosting.

\subsection{Bayesian hyperparameter optimization}
The final optimization stage used \texttt{BayesSearchCV} with 32 iterations over the following random-forest search space:
\begin{itemize}
    \item \texttt{n\_estimators}: 100 to 450,
    \item \texttt{max\_depth}: 8 to 30,
    \item \texttt{min\_samples\_split}: 2 to 12,
    \item \texttt{min\_samples\_leaf}: 1 to 6,
    \item \texttt{max\_features}: \{\texttt{sqrt}, \texttt{log2}, \texttt{None}\},
    \item \texttt{bootstrap}: \{\texttt{True}, \texttt{False}\}.
\end{itemize}
The objective was mean cross-validated $R^2$ over the same stratified 5-fold splits used in the full-dataset evaluation.

\section{Results}

\subsection{Phase I: preliminary model comparison on the 1000-sample dataset}
The first meaningful benchmarking phase used the 1000-sample correlation-preserving synthetic dataset with the original 3 input variables. Table~\ref{tab:phase1} summarizes the optimized model results reported in the project documentation.

\begin{table}[H]
\centering
\caption{Optimized model performance on the 1000-sample improved synthetic dataset using an 80/20 random split.}
\label{tab:phase1}
\begin{tabular}{lcccc}
\toprule
Model & $R^2$ & RMSE & MAE & MSE \\
\midrule
Random Forest & 0.8082 & 0.0409 & 0.0273 & 0.00167 \\
BPNN & 0.7984 & 0.0419 & 0.0311 & 0.00176 \\
KNN & 0.7972 & 0.0420 & 0.0279 & 0.00176 \\
Linear Regression & 0.7398 & 0.0476 & 0.0399 & 0.00227 \\
\bottomrule
\end{tabular}
\end{table}

At this stage, random forest emerged as the strongest model family, though the gap to BPNN and KNN was not yet overwhelming. This suggested that the learning task was nonlinear but not yet dominated by one specific inductive bias under random splitting.

\subsection{Phase II: sequential holdout validation exposed distribution shift}
The project then moved to the harder 3000-sample feature-expanded dataset. A sequential split was adopted to determine whether models trained on the first 2000 samples generalized to the last 1000. The baseline results were catastrophic: BPNN reached $R^2 = -331.96$, linear regression $R^2 = -220.36$, KNN $R^2 = -2.23$, and random forest $R^2 = -1.221$. All models performed worse than predicting the mean bonded thickness.

Diagnostic analysis showed why. The training and validation partitions differed sharply in several variables. For example, concentration increased from a mean of 1.78 g/mL in the training partition to 10.68 g/mL in the validation partition, while total thickness decreased from 1.996 nm to 0.693 nm. Surface tension and viscosity also shifted upward. The validation set therefore represented a substantially different operating regime than the training set.

\subsection{Phase II optimization: feature expansion and robust model selection}
Once the failure mode was identified, the project expanded the feature set from 3 to all 7 available descriptors and focused on tree-based models. The optimized holdout random forest used 200 trees, depth 20, minimum split size 5, minimum leaf size 2, and square-root feature subsampling. Under the same sequential holdout split, performance recovered to
\begin{equation}
R^2 = 0.3284, \quad \mathrm{RMSE} = 0.1065, \quad \mathrm{MAE} = 0.0745.
\end{equation}

This was a major improvement over the baseline holdout random forest ($R^2 = -1.221$), but it still reflected a hostile validation scenario rather than the more typical independent-and-identically-distributed (i.i.d.) generalization setting.

\subsection{Phase III: stratified full-dataset cross-validation}
The next step removed sequential partitioning artifacts by evaluating models with stratified 5-fold cross-validation over the full 3000-sample dataset. Under this design, the optimized random forest achieved:
\begin{equation}
R^2_{\mathrm{OOF}} = 0.971620, \quad \mathrm{RMSE}_{\mathrm{OOF}} = 0.030851, \quad \mathrm{MAE}_{\mathrm{OOF}} = 0.019663.
\end{equation}

The mean fold-level metrics were similarly strong:
\begin{equation}
R^2_{\mathrm{mean}} = 0.971593 \pm 0.003533.
\end{equation}

This result completely changed the project interpretation. The earlier $R^2 \approx 0.33$ figure was not a realistic ceiling for the data-generating process; it was largely a consequence of sequential validation under regime shift. Once the evaluation design matched the full dataset more fairly, the problem became highly learnable.

\subsection{Phase IV: ensemble comparison}
Table~\ref{tab:ensemble} compares the base models and ensembles using out-of-fold performance on the same stratified 5-fold design.

\begin{table}[H]
\centering
\caption{Full-dataset stratified 5-fold out-of-fold performance for base models and ensembles.}
\label{tab:ensemble}
\begin{tabular}{lcccc}
\toprule
Model & $R^2_{\mathrm{OOF}}$ & RMSE$_{\mathrm{OOF}}$ & MAE$_{\mathrm{OOF}}$ & MSE$_{\mathrm{OOF}}$ \\
\midrule
Stacked Ridge (RF + GB) & 0.972202 & 0.030533 & 0.019859 & 0.0009323 \\
Adaptive Blend (RF + GB) & 0.971985 & 0.030652 & 0.019823 & 0.0009395 \\
Fixed Blend 50/50 & 0.971920 & 0.030687 & 0.020147 & 0.0009417 \\
Random Forest & 0.971620 & 0.030851 & 0.019663 & 0.0009518 \\
Gradient Boosting & 0.969843 & 0.031802 & 0.021518 & 0.0010114 \\
\bottomrule
\end{tabular}
\end{table}

The stacked ensemble produced the best ensemble result, though the gain over the standalone random forest was modest: approximately $+0.00058$ in out-of-fold $R^2$. The adaptive blend consistently favored random forest, with inner cross-validation selecting random-forest weights between 0.6 and 0.8 across folds. The stacked ridge meta-learner also assigned slightly larger coefficients to random-forest predictions (approximately 0.53--0.56) than to gradient-boosting predictions (approximately 0.45--0.48), confirming that gradient boosting added information but did not dominate the prediction task.

\subsection{Phase V: Bayesian optimization of the best model}
Because random forest remained the strongest single model and the dominant contributor inside the ensembles, it was chosen for final hyperparameter optimization. Bayesian optimization over 32 iterations produced the following best configuration:
\begin{itemize}
    \item \texttt{bootstrap = True}
    \item \texttt{max\_depth = 30}
    \item \texttt{max\_features = log2}
    \item \texttt{min\_samples\_leaf = 1}
    \item \texttt{min\_samples\_split = 2}
    \item \texttt{n\_estimators = 323}
\end{itemize}

The final tuned random forest achieved:
\begin{equation}
R^2_{\mathrm{OOF}} = 0.973424, \quad \mathrm{RMSE}_{\mathrm{OOF}} = 0.029855, \quad \mathrm{MAE}_{\mathrm{OOF}} = 0.018591.
\end{equation}
Its best cross-validated search score was $R^2 = 0.973392$, and when refit to all 3000 samples it achieved an in-sample $R^2 = 0.996519$.

This tuned random forest exceeded the stacked ensemble and therefore became the overall best model in the project.

\section{Detailed Analysis and Discussion}

\subsection{Why the project initially failed}
The most important lesson from the project is that poor results were initially caused less by ``bad machine learning'' than by a mismatch between data design and evaluation design. The original sequential holdout study conflated two distinct challenges:
\begin{enumerate}
    \item learning the bonded-thickness mapping itself, and
    \item extrapolating into a shifted operating regime.
\end{enumerate}
When those challenges were mixed together, nearly every baseline model collapsed. That collapse was informative: it showed that conventional random train/test splits would not have exposed the regime sensitivity nearly as clearly.

\subsection{Feature completeness mattered more than early model complexity}
The original 3-feature setup included concentration, uncoated layer, and total thickness. Concentration carried meaningful signal, but the other two variables were relatively weak on their own. Once polarity, viscosity, boiling point, and surface tension were added, performance improved sharply under the difficult holdout regime. This shows that the central problem was not only nonlinear response but also incomplete state description.

\subsection{Why random forests were consistently strong}
Random forests performed best in almost every phase for several reasons:
\begin{itemize}
    \item they capture nonlinear interactions without requiring manual feature crossing,
    \item they are robust to monotonic rescaling,
    \item they tolerate mixed variable types and mixed correlation strengths,
    \item they average across trees and thereby reduce variance.
\end{itemize}
These advantages are especially relevant in small-data materials workflows where engineered features are limited and distributional peculiarities are common.

\subsection{Interpreting the ensemble results}
The ensemble study is still valuable even though the final tuned random forest outperformed it. First, the ensemble confirms that gradient boosting contributes partially complementary information. The adaptive blend and stacked ridge both improved on the untuned random forest. Second, the small gain indicates that most of the recoverable signal was already captured by random forest. In other words, the ensemble was useful not because it revolutionized performance, but because it quantified the remaining headroom before full single-model optimization.

\subsection{Final model interpretation through feature importance}
The Bayesian-optimized random forest assigned the following normalized importance weights:
\begin{enumerate}
    \item Concentration: 0.3403
    \item Total Thickness: 0.2895
    \item Uncoated Layer: 0.1751
    \item Polarity: 0.0707
    \item Surface Tension: 0.0489
    \item Viscosity: 0.0466
    \item Boiling Point: 0.0288
\end{enumerate}

This ranking is chemically and process-wise plausible. Concentration is naturally tied to solids content and material loading. Total thickness and uncoated layer encode state information that likely summarizes deposition history or substrate-film interaction. The physicochemical descriptors are individually weaker but still useful as secondary modulators.

\subsection{Residual diagnostics and feature-engineering opportunities}
Residual diagnostics from the best full-dataset ensemble, and by extension the full-dataset modeling phase more broadly, suggest several concrete directions for future feature engineering:
\begin{itemize}
    \item \textbf{Polarity}: absolute residual correlation of approximately 0.212 and a high/low quartile MAE ratio of 1.49 suggest that the model error varies with polarity regime. This points toward interaction terms involving polarity with other solvent descriptors.
    \item \textbf{Surface tension}: nontrivial residual correlation suggests that capillary effects are not fully linearized by the current feature representation. Pairwise interactions with viscosity or concentration may help.
    \item \textbf{Boiling point}: residual structure suggests that a piecewise or spline transformation may be more appropriate than a purely implicit tree representation.
    \item \textbf{Total thickness}: residual dependence and strong positive skew suggest that a log-transformed or quantile-binned version of total thickness may improve robustness.
\end{itemize}

These diagnostics are important because they convert a high-performing black-box result into a more actionable modeling roadmap.

\subsection{Reconciling the holdout and cross-validation narratives}
At first glance, the project's reported performances appear inconsistent: why did the optimized holdout random forest achieve only $R^2 = 0.3284$ while later cross-validated models exceeded $R^2 = 0.97$? The answer is that the two experiments measure different things.

The sequential holdout test asked whether a model trained in one regime could extrapolate into a later regime with materially shifted concentration, thickness, and surface-tension distributions. The stratified cross-validation study instead asked how well the model generalizes under balanced reuse of the full data distribution. Both are legitimate questions, but they should not be interpreted as the same benchmark. In practical terms:
\begin{itemize}
    \item the sequential holdout result measures regime-shift robustness,
    \item the stratified cross-validation result measures i.i.d.-style predictive fidelity across the full dataset.
\end{itemize}
This distinction is essential for any research paper built from the project.

\subsection{Limitations}
Several limitations remain:
\begin{enumerate}
    \item \textbf{Synthetic augmentation dominates the final sample count.} Although the project began from real experimental data, most of the final training pool is synthetic. This is acceptable for algorithm development but weaker than a large purely experimental campaign.
    \item \textbf{Some synthetic values are not strictly physical.} The improved generator allowed 10\% extrapolation beyond original bounds, producing small negative values in concentration and total thickness. Future generators should enforce positivity and chemistry-aware constraints.
    \item \textbf{The final target is empirical bonded thickness, not a direct mechanistic state variable.} This is appropriate for prediction but complicates mechanistic interpretation.
    \item \textbf{External validation is still needed.} The full-dataset results are strong, but additional experimental or process-distinct datasets are needed before deployment claims can be made confidently.
\end{enumerate}

\section{Conclusion}
This paper documented the complete evolution of a machine-learning study for dip-coated thin-film thickness prediction. The project began with a 58-sample experimental dataset, expanded through synthetic augmentation, and matured into a 3000-sample feature-rich study that combined process variables with physicochemical descriptors. Along the way, several key conclusions emerged.

First, classical dip-coating theory remains the correct physical starting point, but it is not sufficient on its own for predicting bonded thickness in a multivariable materials system. Second, model performance depends critically on validation protocol. A sequential holdout split revealed severe distribution shift and drove all baseline models to negative $R^2$, but that failure was diagnostically useful rather than merely discouraging. Third, feature expansion from 3 to 7 variables and a shift toward tree-based methods transformed the holdout random forest from $R^2 = -1.221$ to $R^2 = 0.3284$ under regime-shift validation. Fourth, once the full dataset was evaluated using stratified 5-fold cross-validation, the problem proved highly learnable: the optimized random forest reached $R^2 = 0.9716$, ensembles reached approximately $R^2 = 0.9722$, and Bayesian hyperparameter optimization pushed the final tuned random forest to $R^2 = 0.9734$ out of fold.

The final best model in the project is therefore the Bayesian-optimized random forest, not the stacked ensemble. Its success confirms that the response surface is strongly nonlinear, interaction-rich, and especially well matched to tree ensembles. At the same time, residual diagnostics show that there is still meaningful work to do in feature engineering, particularly around polarity, surface tension, boiling point, and thickness transformations.

More broadly, the project illustrates a useful pattern for small-data materials informatics: begin with physics, augment data cautiously, stress-test validation design, prefer robust nonlinear models when features are heterogeneous, and reserve expensive global hyperparameter search for the strongest model family after the validation framework has stabilized. That workflow turned a project that initially appeared unstable into one with a credible, high-performing predictive ceiling.

\section*{References}
\begin{thebibliography}{99}

\bibitem{landau1942}
L. Landau and B. Levich,
``Dragging of a liquid by a moving plate,''
\textit{Acta Physicochimica U.R.S.S.}, vol. 17, no. 1--2, pp. 42--54, 1942.

\bibitem{rio2017}
E. Rio and F. Boulogne,
``Withdrawing a solid from a bath: How much liquid is coated?''
\textit{Advances in Colloid and Interface Science}, vol. 247, pp. 100--114, 2017.

\bibitem{schmidt2019}
J. Schmidt, M. R. G. Marques, S. Botti, and M. A. L. Marques,
``Recent advances and applications of machine learning in solid-state materials science,''
\textit{npj Computational Materials}, vol. 5, article 83, 2019.

\bibitem{kailkhura2019}
B. Kailkhura, B. Gallagher, S. Kim, A. Hiszpanski, and T. Y.-J. Han,
``Reliable and explainable machine-learning methods for accelerated material discovery,''
\textit{npj Computational Materials}, vol. 5, article 108, 2019.

\bibitem{aha1991}
D. W. Aha, D. Kibler, and M. K. Albert,
``Instance-based learning algorithms,''
\textit{Machine Learning}, vol. 6, no. 1, pp. 37--66, 1991.

\bibitem{rumelhart1986}
D. E. Rumelhart, G. E. Hinton, and R. J. Williams,
``Learning representations by back-propagating errors,''
\textit{Nature}, vol. 323, pp. 533--536, 1986.

\bibitem{breiman2001}
L. Breiman,
``Random forests,''
\textit{Machine Learning}, vol. 45, no. 1, pp. 5--32, 2001.

\bibitem{friedman2001}
J. H. Friedman,
``Greedy function approximation: A gradient boosting machine,''
\textit{The Annals of Statistics}, vol. 29, no. 5, pp. 1189--1232, 2001.

\bibitem{xu2019}
L. Xu, M. Skoularidou, A. Cuesta-Infante, and K. Veeramachaneni,
``Modeling tabular data using conditional GAN,''
arXiv:1907.00503, 2019.

\bibitem{snoek2012}
J. Snoek, H. Larochelle, and R. P. Adams,
``Practical Bayesian optimization of machine learning algorithms,''
in \textit{Advances in Neural Information Processing Systems 25}, 2012, pp. 2951--2959.

\end{thebibliography}

\end{document}
